\begin{tikzpicture}[
    % Ustawienia odstępów
    node distance=1.2cm and 1.0cm,
    auto,
    % Styl bloków
    block/.style={
        rectangle,
        draw=black,
        thick,
        fill=white,
        text width=2.8cm,
        text centered,
        rounded corners,
        minimum height=2cm,
        font=\footnotesize\sffamily
    },
    % Styl linii
    line/.style={
        draw,
        -latex,
        thick,
        gray!80
    },
    % Styl etykiet
    label_text/.style={
        midway,
        font=\scriptsize\itshape,
        align=center,
        color=black!70,
        fill=white,      % Dodano tło dla czytelności, jeśli przetnie linię
        inner sep=2pt
    }
]

    % --- RZĄD 1 (Góra) ---
    \node [block, fill=orange!10] (gitlab) {
        \textbf{1. GitLab}\\
        \textit{(Commit State)}\\[0.3em]
        Repozytorium kodu
    };

    \node [block, right=of gitlab, fill=blue!5] (infra) {
        \textbf{2. OpenTofu}\\
        \textbf{\& Proxmox}\\
        \textit{(HCL State)}\\[0.3em]
        Wdrażanie VM
    };

    \node [block, right=of infra, fill=yellow!10] (k3s) {
        \textbf{3. Ansible \& K3s}\\
        \textit{(Cluster Setup)}\\[0.3em]
        Instalacja K8s
    };

    % --- RZĄD 2 (Dół) ---
    \node [block, below=of k3s, fill=cyan!5] (flux_boot) {
        \textbf{4. Flux Bootstrap}\\
        \textit{(Deploy Operators)}\\[0.3em]
        Inicjalizacja GitOps
    };

    \node [block, left=of flux_boot, fill=cyan!10] (flux_sync) {
        \textbf{5. Flux Reconciliation}\\
        \textit{(Helm \& Kustomize)}\\[0.3em]
        Synchronizacja aplikacji
    };

    % --- POŁĄCZENIA ---
    \path [line] (gitlab) -- (infra);
    \path [line] (infra) -- (k3s);
    \path [line] (k3s) -- (flux_boot);
    \path [line] (flux_boot) -- (flux_sync);
    
    \path [line, dashed, rounded corners] (flux_sync.west) -| node [label_text, pos=0.3] {Monitorowanie zmian} (gitlab.south);

\end{tikzpicture}
